\section{The increasing inverse of Barnes' \texorpdfstring{$G$}{G}-function}
\label{sec:barnes-inverse}

The logarithmic expansion \eqref{eq:barnes-logarithm} permits inversion
in a scale whose dominant term is expressed exactly by the Lambert
function. Keeping that term exact separates the successive powers of
the large source from the slower logarithmic coefficients.

\subsection{The real branch above three}

The Barnes function is strictly increasing on $(3,\infty)$ and maps
this interval onto $(1,\infty)$. To establish monotonicity, put
$v=x-1$ and $C=\log(2\pi)$. Differentiating
\cite[(5.17.4)]{dlmf} gives the exact identity
\begin{equation}
 \frac{d}{dv}\log G(v+1)
   =v\bigl(\psi(v+1)-1\bigr)+\frac{C-1}{2},
 \label{eq:barnes-logarithmic-derivative}
\end{equation}
where $\psi$ is the digamma function. Its defining series gives
\[
 \psi(v+1)=\lim_{n\to\infty}
   \left(\log n-\sum_{k=1}^{n}\frac1{v+k}\right).
\]
The function $t\mapsto1/(v+t)$ is strictly convex. Comparing each
midpoint value with its integral over $[k-1/2,k+1/2]$ gives a positive
gap; the gap from the first interval persists in the limit. Thus
$\psi(v+1)>\log(v+1/2)$ for $v>0$. Consequently the derivative
in \eqref{eq:barnes-logarithmic-derivative} exceeds
\[
 h(v)=v\bigl(\log(v+1/2)-1\bigr)+\frac{C-1}{2}.
\]
For $v\ge2$, one has
$h'(v)=\log(v+1/2)-1/(2v+1)>0$ and $h(2)>0$.
For an elementary strict bound, the inequality
$\log u>2(u-1)/(u+1)$ for $u>1$, applied to the factors in
$5/2=(3/2)(5/3)$, gives
\[
 \log(5/2)=\log(3/2)+\log(5/3)
   >\frac25+\frac12=\frac9{10}>\frac89.
\]
Likewise, $2\pi>6=(3/2)2^2$ gives
\[
 C>\log(3/2)+2\log2
   >\frac25+\frac43=\frac{26}{15}>\frac53.
\]
These bounds imply $h(2)>1/9$.
Thus the derivative is positive for $x>3$.
The recurrence gives $G(3)=1$, and \eqref{eq:barnes-logarithm}
gives $G(x)\to\infty$. Denote this increasing inverse by $g(y)$.
All subsequent assertions concern $y\to+\infty$ on this branch.

\subsection{An exact Lambert balance}

Write $Y=\log y$ and define
\begin{equation}
 X=\sqrt{\frac{4Y}{W_0(4Y/e^3)}},\qquad
 L=\log X,\qquad t=\frac1X,\qquad q=\frac1{L-1},
 \label{eq:barnes-inverse-coordinate}
\end{equation}
where $W_0$ is the principal real Lambert function. Then
\begin{equation}
 \frac{X^2}{2}\left(\log X-\frac32\right)=Y,
 \qquad
 L-1=\frac{1+W_0(4Y/e^3)}2>0.
 \label{eq:barnes-inverse-balance}
\end{equation}
As $y\to\infty$, both $X$ and $L$ tend to infinity. The leading
balance and monotonicity imply $(g(y)-1)/X\to1$.
The coefficient coordinate $q$ uses $L-1$ because the derivative
of the dominant logarithmic phase is $X(L-1)$.

Let $A=\log\mathcal A$, where $\mathcal A$ is Glaisher's constant,
so that $\zeta'(-1)=1/12-A$. Write
$g(y)=1+X(1+U)$. After substituting the finite expansion
\eqref{eq:barnes-logarithm} and dividing the logarithmic residual
by $X^2(L-1)$, the normalized finite equation is
$\Psi_m(U,t,q)=0$, where
\begin{equation}
\begin{aligned}
 \Psi_m(U,t,q)={}&U+\frac{U^2}{2}
   +\frac q2\left((1+U)^2\log(1+U)-U-\frac{U^2}{2}\right)\\
 &+\frac{Cqt}{2}(1+U)-t^2\left(\frac1{12}+Aq\right)
   -\frac{qt^2}{12}\log(1+U)\\
 &+q\sum_{k=1}^{m}
   \frac{B_{2k+2}\,t^{2k+2}(1+U)^{-2k}}{2k(2k+2)}.
\end{aligned}
 \label{eq:barnes-inverse-normalized-equation}
\end{equation}
The omitted part of the true normalized equation is
$\OO(qt^{2m+4})$, uniformly for $U$ in a fixed sufficiently small
real interval. The finite equation is analytic near the origin in
the three independent variables $U,t,q$ and has derivative one
with respect to $U$ there.

\subsection{Polynomial coefficients and three complete blocks}

Temporarily regard $t$ and $q$ as independent. Seek a formal solution
\begin{equation}
 U=\sum_{n\ge0}b_n(q)t^{n+1},\qquad
 b_n\in\mathbb Q[A,C,q].
 \label{eq:barnes-inverse-unit-series}
\end{equation}
With $U_{<n}=\sum_{j=0}^{n-1}b_j(q)t^{j+1}$, the triangular recurrence is
\begin{equation}
 b_n(q)=-[t^{n+1}]\Psi_m(U_{<n},t,q),\qquad
 m\ge\max\left(0,\left\lfloor\frac{n-1}{2}\right\rfloor\right).
 \label{eq:barnes-inverse-coefficient-recurrence}
\end{equation}
After the linear terms inside the logarithmic bracket cancel,
only the first term $U$ contains the new coefficient at the requested
power; every other occurrence has an additional factor $t$ or $U$.
This proves uniqueness and polynomiality by induction. A larger $m$
does not change coefficients already determined.

The first coefficients are
\begin{align}
 b_0(q)&=-\frac{Cq}{2},\notag\\
 b_1(q)&=\frac1{12}+Aq+\frac{C^2q^2}{8}-\frac{C^2q^3}{8},\notag\\
 b_2(q)&=\frac{CAq^3}{2}+\frac{C^3q^4}{12}-\frac{C^3q^5}{16}.
 \label{eq:barnes-inverse-first-coefficients}
\end{align}
Thus the first three complete blocks of the original source are
\begin{equation}
 G_3(y)=X+\left(1-\frac{Cq}{2}\right)
       +\frac1X\left(\frac1{12}+Aq
                    +\frac{C^2q^2}{8}-\frac{C^2q^3}{8}\right).
 \label{eq:barnes-inverse-three-blocks}
\end{equation}
The source translation by one belongs to the block at power $X^0$.
Each block retains its whole polynomial in $q$. Since
$X^{-1}(L-1)^d\to0$ for every fixed real $d$, distinct nonzero
blocks are ordered by their powers of $X^{-1}$ even when their
coefficient polynomials have different orders of vanishing at $q=0$.

\begin{theorem}[Finite-order inverse Barnes expansion]
\label{thm:barnes-inverse-finite-order}
For every fixed integer $N\ge0$,
\begin{equation}
 g(y)=1+X+\sum_{n=0}^{N}X^{-n}b_n(q)+\OO(X^{-N-1}).
 \label{eq:barnes-inverse-finite-remainder}
\end{equation}
In particular,
\begin{equation}
 g(y)-G_3(y)=\frac{b_2(q)}{X^2}+\OO(X^{-3})
            =\OO\left(\frac1{X^2(L-1)^3}\right),
 \label{eq:barnes-inverse-three-remainder}
\end{equation}
and
\begin{equation}
 \lim_{y\to\infty}X^2(L-1)^3\bigl(g(y)-G_3(y)\bigr)
   =\frac{\log(2\pi)\log\mathcal A}{2}>0.
 \label{eq:barnes-inverse-three-error-limit}
\end{equation}
\end{theorem}
\begin{proof}
Fix $N$ and choose
$m\ge\max(0,\lfloor(N-1)/2\rfloor)$. The analytic implicit function
theorem gives a solution $U_m(t,q)$ of \eqref{eq:barnes-inverse-normalized-equation}
in a common neighborhood of $(t,q)=(0,0)$. Taylor expansion in $t$,
uniformly for $q$ in a smaller neighborhood, gives
\[
 U_m(t,q)=\sum_{n=0}^{N}b_n(q)t^{n+1}+\OO(t^{N+2}).
\]
For the true normalized equation, differentiation with respect to $U$
gives
\[
 \frac{\left.\dfrac{d}{dv}\log G(v+1)\right|_{v=X(1+U)}}{X(L-1)}.
\]
By \eqref{eq:barnes-logarithmic-derivative} and the digamma asymptotics,
this tends uniformly to $1+U$ on any fixed sufficiently small interval
about zero. It is therefore bounded away from zero there.
The true inverse has $U\to0$. The mean value theorem and the finite
Barnes remainder bound its difference from $U_m$ by
$\OO(qt^{2m+4})$. Multiplication by $X=t^{-1}$ and the inequality
$2m+3\ge N+1$ prove \eqref{eq:barnes-inverse-finite-remainder}.
For \eqref{eq:barnes-inverse-three-remainder}, use $N=2$ and separate
the $b_2$ term. Its leading coefficient is $CAq^3/2$, and
$(L-1)^3/X\to0$, which also proves
\eqref{eq:barnes-inverse-three-error-limit}.
\end{proof}

The generic finite-order remainder has no obligatory factor $q$.
Indeed, at $q=0$ the analytic finite-model solution is
$U_m(t,0)=-1+\sqrt{1+t^2/6}$, which already has nonzero coefficients
at arbitrarily high even powers of $t$. The proof concerns each fixed
finite model and makes no assertion that the inverse series converges
as its truncation order tends to infinity.

\subsection{Affine reconstruction and real powers}

Let $H(Y)$ be the increasing solution of $\log G(H(Y))=Y$, so
$g(y)=H(\log y)$. For fixed real parameters
$\alpha a r\ne0$, the equation
\[
 y=d+a\,G(\alpha x+\beta)^r,\qquad \alpha x+\beta>3,
\]
has the exact source reconstruction
\begin{equation}
 x=\frac{H(Y)-\beta}{\alpha},\qquad
 Y=\frac1r\log\frac{y-d}{a}>0,
 \qquad \frac{y-d}{a}>0.
 \label{eq:barnes-inverse-affine-reconstruction}
\end{equation}
The regime $Y\to\infty$ corresponds to $(y-d)/a\to\infty$ for
$r>0$, and to $(y-d)/a\to0^+$ for $r<0$. The latter approaches
the finite target $d$ from the side determined by $a$.
An affine logarithmic target $y=d+a\log G(\alpha x+\beta)$ instead
uses $Y=(y-d)/a$. All these assertions keep the parameters fixed.

A fixed real power $x^\nu$ is defined eventually for $\alpha>0$;
for $\alpha<0$, restrict here to integer $\nu$. Its normalized
expression is
\[
 x^\nu=\alpha^{-\nu}X^\nu
       \bigl(1+(1-\beta)t+U\bigr)^\nu.
\]
Finite analytic composition therefore gives polynomial coefficients
in $q$ at successive powers of $X^{-1}$, with the same argument as
\eqref{eq:gamma-inverse-observable-expansion}. If the first omitted
nonzero coefficient at $X^{\nu-j}$ is
$\gamma q^h+\OO(q^{h+1})$, then the corresponding error is
asymptotic to $\gamma X^{\nu-j}(L-1)^{-h}$: the next full core power
has size at most $\OO(X^{\nu-j-1})$, and
$X^{-1}(L-1)^h\to0$. Entire zero blocks are omitted before selecting
this first coefficient.
